\documentclass[10pt,a4paper]{amsart}
\usepackage[margin=19mm]{geometry}
\usepackage{amsmath,amssymb,mathtools}
\usepackage[scr=rsfs,cal=euler]{mathalfa}
\usepackage{xfrac}
\usepackage[unicode,hidelinks,bookmarks=false]{hyperref}
\newcommand{\ie}{i.e.\ }
\newcommand{\cf}{cf.\ }
\newcommand{\resp}{respectively\ }
\newcommand{\Subsection}{\subsection}
\providecommand{\nobreakdash}{\nobreak\hbox{-}}
\setlength{\emergencystretch}{2em}
\multlinegap0pt
\usepackage{bm}
\usepackage{bbm}
\usepackage{dsfont}
\usepackage{xspace}
\usepackage{cancel}
\usepackage{mathtools}
\usepackage{amsthm}
\makeatletter
\@ifundefined{lem}{\newtheorem{lem}{Lemma}}{}
\makeatother
\title[Numbers with Four Close Factorizations]{Numbers with Four Close Factorizations}
\author{Tsz Ho Chan, Laura Holmes, Michael Liu, Jose Villarreal}
\date{Source: arXiv:2508.02818v1 (CC BY 4.0)}
\begin{document}
\maketitle

\noindent\emph{Source and licence.} Numbers with Four Close Factorizations. Version arXiv:2508.02818v1 (CC BY 4.0), distributed under the Creative Commons Attribution 4.0 International licence, as recorded in the arXiv licence metadata for this exact version and shown on the abstract page https://arxiv.org/abs/2508.02818v1. Full rights notice: licenses/arxiv-2508.02818-CC-BY-4.0.txt.\par\smallskip

\noindent\textit{Editorial scope.} Counted unit: the Lem whose source label is \texttt{lem14} in the article (source0.tex lines 912--918, source label \texttt{lem14}), delivered together with the article's own complete original proof of it (source0.tex lines 920--954, one \texttt{proof} environment of 35 lines). No mathematics is written, repaired or re-derived: statement and proof are the article's own text line for line. Required context retained uncounted: basic3' (lines 156--158); lem4 (lines 194--196); lem5 (lines 201--203); definitions (lines 457--466); km-g (lines 502--504); a3b3 (lines 507--509); alpha2beta2 (lines 519--521); pell-1-g (lines 523--525). Every definition, display and label of those lines is retained unchanged. Omitted from this package and not redistributed: 1--155, 159--193, 197--200, 204--456, 467--501, 505--506, 510--518, 522--522, 526--911, 919--919, 955--1074. Nothing inside a retained excerpt is omitted or altered: every retained body is delivered exactly as the article's lines are written, so no omission receipt of any kind accompanies this package. Citations: the retained excerpts carry no citation command at all, so no bibliography entry of the article is transcribed and no reference is redistributed. Reference adaptation: none was needed and none was made; every reference inside the delivered unit resolves inside it, so no unresolved reference or citation is produced. Printed slips kept literal: no printed word, symbol, hypothesis or number of the counted statement or of its proof was altered. Layout and class: the article's own class is \texttt{article} with packages including graphicx, amsmath, amssymb, amsthm, xcolor, tabularx, appendix, multirow; the wrapper is the lane's pinned \texttt{amsart} editorial wrapper at 10pt on A4, which restates the theorem-like environments the retained text needs and adds the article's own macro definitions that the retained text needs (none), with extra wrapper packages (bm, bbm, dsfont, xspace, cancel, mathtools). The delivered numbering is the unit's own. Assets: no retained span contains a figure environment, a graphics-inclusion command, a PDF-page inclusion, a table or a TikZ picture; no image file is copied into this package, the package declares no asset dependency and no image file is redistributed with it. Licence: version arXiv:2508.02818v1 of this article is distributed under the Creative Commons Attribution 4.0 International licence, as recorded in the arXiv licence metadata for this exact version and shown on the abstract page https://arxiv.org/abs/2508.02818v1. A full rights notice is delivered at licenses/arxiv-2508.02818-CC-BY-4.0.txt. Characters: every retained excerpt is pure ASCII, so the delivered engine, XeLaTeX with its default Latin Modern text font, reports no missing character.
\begin{equation} \label{basic3'}
A = \frac{a_2 a_1 (b_2 - b_1)}{D_{21}} = \frac{a_3 a_1 (b_3 - b_1)}{D_{31}} = \frac{a_3 a_2 (b_3 - b_2)}{D_{32}}.
\end{equation}

\begin{lem} \label{lem4}
We have the inequality $D_{31} > D_{21}$.
\end{lem}

\begin{lem} \label{lem5}
We have the inequality $D_{32} + D_{21} > D_{31}$.    
\end{lem}

\begin{equation} \label{definitions}
d_a := \gcd(a_1, a_2), \; d_b := \gcd(b_1, b_2), \; \left\{ \begin{array}{l}
a_1 := d_a \alpha_1 \\
a_2 := d_a \alpha_2,
\end{array} \right. \; 
\left\{ \begin{array}{l}
b_1 := d_b \beta_1 \\
b_2 := d_b \beta_2,
\end{array} \right.
\end{equation}

\begin{equation} \label{km-g}
k m = \frac{d_a d_b}{D_{21}}{D_{32} D_{31} (D_{32} +D_{21} - D_{31})} .
\end{equation}

\begin{equation} \label{a3b3}
a_3 = \frac{d_a (D_{31} \alpha_2 - D_{32} \alpha_1)}{D_{21}} \; \; \text{ and } \; \; b_3 = \frac{d_b (D_{31} \beta_2 - D_{32} \beta_1)}{D_{21}}.
\end{equation}

\begin{equation} \label{alpha2beta2}
\alpha_2 = \frac{D_{21} k \beta_1 + d_a D_{31} D_{32} \alpha_1}{d_a D_{31} (D_{31} - D_{21})} \; \; \text{ and } \; \; \beta_2 = \frac{D_{21} m \alpha_1 + d_b D_{31} D_{32} \beta_1}{d_b D_{31} (D_{31} - D_{21})}.
\end{equation}

\begin{equation} \label{pell-1-g}
d_b k \beta_1^2 - d_a m \alpha_1^2 = D_{31} (D_{31} - D_{21}).
\end{equation}

\begin{lem} \label{lem14}
For $1 \le D_{21}, D_{31}, D_{32} \le 5$, we have
\begin{equation*}
    R_4 = \frac{m (D_{31}-D_{21}) \bigl( D_{21}\sqrt{\frac{k m}{d_a d_b}}+D_{32}{D_{31} \bigr)\bigl[ D_{21} + D_{31}\sqrt{\frac{d_a d_b}{k m}}(D_{32}-D_{31}+D_{21})}\bigr]}
    {D_{21}D_{31}^2 d_a\bigl( \sqrt{\frac{k m}{d_a d_b}}+D_{32} \bigr)^3} + O\Bigl(\frac{1}{\alpha_1^2}\Bigr).
\end{equation*}
\end{lem}

\begin{proof}
Since $d_a d_b = \gcd(a_1, a_2) \gcd(b_1, b_2) \mid a_2 b_1 - b_2 a_1 = D_{21} \le 5$, we have $1 \le d_a d_b \le 5$. Combining this with \eqref{km-g} and Lemmas \ref{lem4} and \ref{lem5}, we have $1 \le k m \le 5 \cdot 5 \cdot 5 \cdot 4 = 500$. From the Pell-type equation \eqref{pell-1-g}, we obtain
\[
\Bigl( \beta_1 - \sqrt{\frac{d_a m}{d_b k}} \alpha_1 \Bigr) \Bigl( \beta_1 + \sqrt{\frac{d_a m}{d_b k}} \alpha_1 \Bigr) = \beta_1^2 - \frac{d_a m}{d_b k} \alpha_1^2 = O(1)
\]
which implies
\begin{equation} \label{lem10}
\frac{\beta_1}{\alpha_1} - \sqrt{\frac{d_a m}{d_b k}} = O \Bigl( \frac{1}{\sqrt{(d_a m) / (d_b k)} \alpha_1^2} \Bigr) = O \Bigl( \frac{1}{\alpha_1^2} \Bigr)
\end{equation}
as $\frac{d_b k}{d_a m} \le \frac{5 \cdot 500}{1 \cdot 1}$. 

\bigskip

Next, we find formulas for $a_3$, $A$, and $R_4$. Substituting \eqref{alpha2beta2} into \eqref{a3b3}, we get
\begin{equation} \label{a3-new}
a_3 = \frac{k \beta_1 + d_a D_{32} \alpha_1}{D_{31} - D_{21}}.
\end{equation}
By \eqref{basic3'}, \eqref{definitions} and \eqref{a3-new}, we have $A = \frac{a_2 a_1 (b_2 - b_1)}{D_{21}} = \frac{d_a^2 d_b \alpha_2 \alpha_1 (\beta_2 - \beta_1)}{D_{21}}$. Hence, we obtain
\begin{equation} \label{R4-temp1}
R_4 = \frac{A}{a_3^3} = \frac{(D_{31} - D_{21})^3 d_a^2 d_b \alpha_2 \alpha_1 (\beta_2 - \beta_1)}{D_{21} (k \beta_1 + d_a D_{32} \alpha_1)^3}.
\end{equation}
From \eqref{alpha2beta2}, we have $\beta_2 - \beta_1 = \frac{D_{21} m \alpha_1 + d_b D_{31} (D_{32} + D_{21} - D_{31}) \beta_1}{d_b D_{31} (D_{31} - D_{21})}$. Substituting this and \eqref{alpha2beta2} into \eqref{R4-temp1}, we obtain
\begin{align*}
R_4 =& \frac{(D_{31} - D_{21}) d_a \alpha_1 (D_{21} k \beta_1 + d_a D_{31} D_{32} \alpha_1) (D_{21} m \alpha_1 + d_b D_{31} (D_{32} + D_{21} - D_{31}) \beta_1)}{D_{21} D_{31}^2 (k \beta_1 + d_a D_{32} \alpha_1)^3} \\
=& \frac{(D_{31} - D_{21}) d_a \bigl( D_{21} k \frac{\beta_1}{\alpha_1} + d_a D_{31} D_{32} \bigr) \bigl(D_{21} m + d_b D_{31} (D_{32} + D_{21} - D_{31}) \frac{\beta_1}{\alpha_1} \bigr)}{D_{21} D_{31}^2 (k \frac{\beta_1}{\alpha_1} + d_a D_{32})^3}.
\end{align*}
Finally, we apply \eqref{lem10} to the above equation and get
\begin{align*}
R_4 =& \frac{(D_{31} - D_{21}) d_a^2 \bigl( D_{21} \sqrt{\frac{k m}{d_a d_b}} + D_{31} D_{32} + O( \frac{1}{\alpha_1^2} ) \bigr)}{D_{21} D_{31}^2 d_a^3 \bigl(\sqrt{\frac{k m}{d_a d_b}} + D_{32} + O( \frac{1}{\alpha_1^2} ) \bigr)^3} \\
&\cdot  \Bigl(D_{21} m + D_{31} (D_{32} + D_{21} - D_{31}) \sqrt{\frac{d_a d_b m}{k}} + O \Bigl( \frac{1}{\alpha_1^2} \Bigr) \Bigr) \\
=& \frac{m (D_{31} - D_{21}) \bigl( D_{21} \sqrt{\frac{k m}{d_a d_b}} + D_{31} D_{32} \bigr) \bigl( 1 + O( \frac{1}{\alpha_1^2} ) \bigr)}{D_{21} D_{31}^2 d_a \bigl(\sqrt{\frac{k m}{d_a d_b}} + D_{32} \bigr)^3 \bigl( 1 + O( \frac{1}{\alpha_1^2} ) \bigr)^3} \\
&\cdot  \Bigl[ D_{21} + D_{31} (D_{32} + D_{21} - D_{31}) \sqrt{\frac{d_a d_b}{k m}} \Bigr] \Bigl(1 + O \Bigl( \frac{1}{\alpha_1^2} \Bigr)  \Bigr).
\end{align*}
This gives the lemma.
\end{proof}

\end{document}
